% ======================================================================
%  sec28.tex —— 第 28 节 （极）丰沛线丛
%  底本：原书 pp.141–156；原书标题：28. (Very) ample line bundles
% ======================================================================

% ---------- 本节用到的宏（若全书导言区已定义则自动跳过） ----------
\providecommand{\Sym}{\operatorname{Sym}}
\providecommand{\QCoh}{\mathsf{QCoh}}
\providecommand{\Sub}{\operatorname{Sub}}

\section{（极）丰沛线丛}\label{sec:28}

\subsection{线丛的主开子集}\label{sec:28.1}

以下结果都可容易地化到 $\cL=\cO_X$ 的情形。

\begin{propdef}\label{prop:28.1.1}
设 $X$ 是概形，$\cL$ 是 $X$ 上的线丛，$f\in\cL(X)$ 是整体截面。对任意 $x\in X$，考虑 $f$ 在 $x$ 处的芽 $f_x\in\cL_x$，并令
\[
X_f := \{x\in X \mid f_x\notin \mathfrak m_x\cL_x\},
\]
其中 $\mathfrak m_x\subseteq\cO_{X,x}$ 是唯一极大理想。则 $X_f$ 是 $X$ 的开子集。称之为对应于 $f$ 的\emph{主开子集}。
\end{propdef}

\begin{yizhu}
原书小节标题中的 principle 是 principal 的误拼，实指「主开子集」。当 $\cL=\cO_X$ 时，$X_f$ 就是 $\Spec R$ 中的标准开集 $D(f)$；一般情形下主开子集对乘法封闭，即 $X_f\cap X_{f'}=X_{f\otimes f'}$，其中 $f'$ 是某个 $\cL'^{\otimes n'}$ 的整体截面。在局部平凡化下，主开子集与仿射标准开集重合，因此丰沛性可以逐个仿射开集检验。
[来源：Stacks Project, Tag 01PR（Section 28.27，Ample invertible sheaves 开篇评注：$X_s$ 开且 $X_s\cap X_{s'}=X_{ss'}$），\url{https://stacks.math.columbia.edu/tag/01PR}（2026-11-17 访问）。]
\end{yizhu}

\begin{lemma}\label{lem:28.1.2}
设 $X$ 是概形，$\cL$ 是 $X$ 上的线丛，$f\in\cL(X)$ 是整体截面。则 $X_f$ 正是 $f\colon\cO_X\to\cL$ 可逆的开区。更精确地说：
\begin{enumerate}[label=(\arabic*),leftmargin=2.4em,itemsep=0.25ex]
  \item 限制 $f|_{X_f}\colon\cO_X|_{X_f}\to\cL|_{X_f}$ 是同构。
  \item 若 $V\subseteq X$ 是子集且 $f|_V$ 是同构，则 $V\subseteq X_f$。
\end{enumerate}
\end{lemma}

\begin{yizhu}
引理 (2) 说明 $X_f$ 是使 $f$ 成为可逆截面的最大开子集，因此 $X_f$ 只由 $f$ 与 $\cL$ 决定，与所选平凡化覆盖无关。另有一条常用事实：对任意仿射开子集 $U\subseteq X$，交 $U\cap X_f$ 仍仿射，故 $X_f$ 的仿射性可以逐块检验。
[来源：Stacks Project, Tag 01PV（Lemma 28.27.4：若 $U\subseteq X$ 仿射开且 $s\in\Gamma(X,\cL)$，则 $U\cap X_s$ 仿射），\url{https://stacks.math.columbia.edu/tag/01PV}（2026-11-17 访问）。]
\end{yizhu}

\begin{exercise}\label{exc:28.1.3}
设 $\cL$ 是 $X$ 上的线丛。证明：典范同构 $p\colon\PP_X(\cL)\to X$ 把 $U_+(f)$ 送到 $X_f$。
\end{exercise}

\begin{construction}\label{con:28.1.4}
设 $X$ 是拟紧且拟分离的概形，$\cL$ 是 $X$ 上的线丛，$f\in\cL(X)$ 是整体截面。对任意拟凝聚 $\cO_X$-模 $\cF$，考虑
\[
M := \bigoplus_{n\in\Z}\Gamma\Bigl(X,\cF\otimes_{\cO_X}\cL^{\otimes n}\Bigr),
\]
它是
\[
S := \bigoplus_{n\ge 0}\Gamma(X,\cL^{\otimes n})
\]
的分次模。考虑限制映射
\[
M \to \bigoplus_{n\in\Z}\Gamma\Bigl(X_f,\cF\otimes_{\cO_X}\cL^{\otimes n}\Bigr).
\]
注意 $S$ 作用在靶上，且上述映射是 $S$-线性的。又由引理 \ref{lem:28.1.2}，$f\in S$ 在靶上可逆地作用。因此有 $S_f$-线性映射
\[
M_f \to \bigoplus_{n\in\Z}\Gamma\Bigl(X_f,\cF\otimes_{\cO_X}\cL^{\otimes n}\Bigr),
\]
特别地有 $S_{(f)}$-线性映射
\begin{equation}
M_{(f)} \to \Gamma(X_f,\cF). \label{eq:28.1}
\end{equation}
\end{construction}

下述结果推广了引理 \ref{lem:20.4.5}。

\begin{theorem}\label{thm:28.1.5}
典范映射 \eqref{eq:28.1} 是同构。
\end{theorem}

\begin{remark}\label{rem:28.1.6}
具体地说，定理 \ref{thm:28.1.5} 断言：
\begin{enumerate}[label=(\arabic*),leftmargin=2.4em,itemsep=0.25ex]
  \item 若 $s\in\Gamma(X,\cF\otimes\cL^{\otimes n})$ 满足 $s|_{X_f}=0$，则当 $m\gg 0$ 时 $s\otimes f^m=0$。
  \item 若 $u\in\Gamma(X_f,\cF)$，则当 $m\gg 0$ 时 $u\otimes (f|_{X_f})^m$ 可延拓为 $X$ 上的截面。
\end{enumerate}
事实上，(1) 等价于 \eqref{eq:28.1} 是单射，而 (2) 等价于 \eqref{eq:28.1} 是满射。
\end{remark}

\begin{proof}
先证 (1)。设 $s\in\Gamma(X,\cF\otimes_{\cO_X}\cL^{\otimes n})$ 满足 $s|_{X_f}=0$。对任意使 $\cL|_U$ 平凡的仿射开子集 $U\subseteq X$，注意 $X_f\cap U=U_{f|_U}$ 是 $U$ 的标准开子集。故 $s|_{X_f\cap U}=0$ 蕴含当 $m\gg 0$ 时
\[
s|_U\otimes (f|_U)^m = 0.
\]
由于 $X$ 拟紧，得到当 $m\gg 0$ 时 $s\otimes f^m=0$。

现证 (2)。设 $u\in\Gamma(X_f,\cF)$。取有限开覆盖 $X=\bigcup_{i\in I}U_i$，其中每个 $U_i$ 仿射且 $\cL|_{U_i}$ 平凡。按 (1) 中同样的论证，当 $m\gg 0$ 时
\[
u\otimes (f|_{X_f\cap U_i})^m
\]
可延拓为截面
\[
u_i \in \Gamma\Bigl(U_i,\cF\otimes_{\cO_X}\cL^{\otimes m}\Bigr),
\]
其中 $i\in I$ 任意。由构造，对 $(i,j)\in I^2$ 与 $U_{ij}:=U_i\cap U_j$，$u_i|_{U_{ij}}-u_j|_{U_{ij}}$ 在 $X_f\cap U_{ij}$ 上的限制为零。故由 (1)（施于概形 $U_{ij}$ 与截面 $f|_{U_{ij}}$），当 $m'\gg 0$ 时
\[
(u_i|_{U_{ij}}-u_j|_{U_{ij}})\otimes (f|_{U_{ij}})^{m'} = 0
\]
对每一对 $(i,j)\in I$ 成立。这说明可以把截面 $u_i\otimes (f|_{U_i})^{m'}$ 粘成整体截面
\[
s\in\Gamma\Bigl(X,\cF\otimes_{\cO_X}\cL^{\otimes (m+m')}\Bigr).
\]
由构造，$s$ 延拓了 $u\otimes (f|_{X_f})^{m+m'}$。
\end{proof}

\subsection{与线丛相伴的态射}\label{sec:28.2}

\begin{construction}\label{con:28.2.1}
设 $f\colon X\to Y$ 是概形态射，$\cL$ 是 $X$ 上的线丛。假设 $f_*(\cL^{\otimes n})$ 对任意 $n\ge 0$ 都拟凝聚。典范态射
\[
\psi\colon \bigoplus_{n\ge 0}f^*f_*(\cL^{\otimes n}) \to \bigoplus_{n\ge 0}\cL^{\otimes n} \simeq \Sym_{\cO_X}(\cL)
\]
诱导态射
\[
\mathfrak r_{\cL/Y}\colon U_{\cL/Y} \to \Proj_Y\Bigl(\bigoplus_{n\ge 0}f_*(\cL^{\otimes n})\Bigr),
\]
其中
\[
U_{\cL/Y} := U_+(\psi) \subseteq \Proj_X(\Sym_{\cO_X}(\cL)) \simeq X
\]
视为 $X$ 的开子概形。

当 $Y=\Spec(\Z)$ 时，有典范态射
\[
\mathfrak r_{\cL}\colon U_{\cL} \to \Proj\Bigl(\bigoplus_{n\ge 0}\cL^{\otimes n}(X)\Bigr).
\]
\end{construction}

\begin{exercise}\label{exc:28.2.2}
证明：对任意截面
\[
s\in\cL^{\otimes n}(X)\simeq\Gamma(Y,f_*(\cL^{\otimes n})),\quad n>0,
\]
有
\[
\mathfrak r_{\cL/Y}^{-1}\bigl(U_+(s)\bigr) = X_s.
\]
特别地，
\[
U_{\cL} = \bigcup_{s\in\cL^{\otimes n}(X),\ n>0}X_s.
\]
\end{exercise}

\begin{construction}\label{con:28.2.3}
设 $f\colon X\to Y$ 是概形态射，$\cL$ 是 $X$ 上的线丛。假设 $f_*(\cL)$ 拟凝聚。典范态射
\[
\psi'\colon \Sym_{\cO_X}(f^*f_*(\cL)) \to \Sym_{\cO_X}(\cL)
\]
诱导态射
\[
\mathfrak s_{\cL/Y}\colon V_{\cL/Y} \to \PP_Y(f_*\cL),
\]
其中
\[
V_{\cL/Y} := U_+(\psi') \subseteq \Proj_X(\Sym_{\cO_X}(\cL)) \simeq X
\]
视为 $X$ 的开子概形。

当 $Y=\Spec(\Z)$ 时，有典范态射
\[
\mathfrak s_{\cL}\colon V_{\cL} \to \PP(\cL(X)).
\]
\end{construction}

\begin{exercise}\label{exc:28.2.4}
证明：对任意截面
\[
s\in\cL(X)\simeq\Gamma(Y,f_*(\cL)),
\]
有
\[
\mathfrak s_{\cL/Y}^{-1}\bigl(U_+(s)\bigr) = X_s.
\]
特别地，
\[
V_{\cL/Y} = \bigcup_{s\in\cL(X)}X_s \subseteq U_{\cL/Y}.
\]
\end{exercise}

\begin{exercise}\label{exc:28.2.5}
证明：$V_{\cL/Y}=X$ 当且仅当典范态射
\[
f^*f_*(\cL) \to \cL
\]
是满射。
\end{exercise}

\subsection{丰沛线丛}\label{sec:28.3}

本小节研究构造 \ref{con:28.2.3} 中的态射 $\mathfrak s_{\cL}$。

\begin{propdef}\label{prop:28.3.1}
设 $X$ 是拟紧概形，$\cL$ 是 $X$ 上的线丛。以下条件等价：
\begin{enumerate}[label=(\alph*),leftmargin=2.6em,itemsep=0.3ex]
  \item 主开子集 $X_f$ 构成 $X$ 的一组拓扑基，其中 $f$ 取遍 $\bigoplus_{n>0}\cL^{\otimes n}(X)$ 中的元素。
  \item[\textup{(a')}] 主开子集 $X_f$ 覆盖 $X$，其中 $f$ 取遍 $\bigoplus_{n>0}\cL^{\otimes n}(X)$ 中使 $X_f$ 仿射的元素。
  \item[\textup{(a'')}] 主开子集 $X_f$ 构成 $X$ 的一组拓扑基，其中 $f$ 取遍 $\bigoplus_{n>0}\cL^{\otimes n}(X)$ 中使 $X_f$ 仿射的元素。
  \item 有 $U_{\cL}=X$，且态射 $\mathfrak r_{\cL}\colon X\to\Proj(\bigoplus_{n>0}\cL^{\otimes n}(X))$ 是像稠密的开浸入。
  \item[\textup{(b')}] 有 $U_{\cL}=X$，且态射 $\mathfrak r_{\cL}\colon X\to\Proj(\bigoplus_{n>0}\cL^{\otimes n}(X))$ 诱导到其像上的同胚。
  \item 对任意拟凝聚 $\cO_X$-模 $\cF$，典范态射
  \[
  \bigoplus_{n\ge N}\cL^{\otimes -n}\otimes_{\Z}\Gamma\Bigl(X,\cF\otimes_{\cO_X}\cL^{\otimes n}\Bigr) \to \cF
  \]
  对任意整数 $N$ 都是满射。
  \item[\textup{(c')}] 对任意拟凝聚理想 $\cI\subseteq\cO_X$，典范态射
  \[
  \bigoplus_{n>0}\cL^{\otimes -n}\otimes_{\Z}\Gamma\Bigl(X,\cI\otimes_{\cO_X}\cL^{\otimes n}\Bigr) \to \cI
  \]
  是满射。
\end{enumerate}
若 $\cL$ 满足上述条件，则称 $\cL$ 是\emph{丰沛的}。
\end{propdef}

\begin{yizhu}
条件 (c) 中的张量积是 $\Z$-模张量积，因为整体截面 $\Gamma(X,\cF\otimes\cL^{\otimes n})$ 只有 $\Z$-模结构，写成 $\otimes_{\cO_X}$ 并无意义。Stacks 把同一条件表述为「各 $n$ 的像之和等于 $\cF$」，且只要求 $n\ge 1$；本书「对任意整数 $N$」的形式是它的加强版，两者等价。
[来源：Stacks Project, Tag 01Q3（Proposition 28.27.13，条件 (5)：$\bigoplus_{n\ge0}\Gamma(X,\cF\otimes\cL^{\otimes n})\otimes_\Z\cL^{\otimes -n}$ 的像之和等于 $\cF$），\url{https://stacks.math.columbia.edu/tag/01Q3}（2026-11-17 访问）。]
\end{yizhu}

\begin{proof}
全证明中，对 $\cO_X$-模 $\cF$ 记
\[
\cF(n) := \cF\otimes_{\cO_X}\cL^{\otimes n}.
\]
又记 $\cO:=\cO_X$ 与
\[
S := \bigoplus\Gamma(X,\cO(n)).
\]
只需证明以下蕴含关系：
\begin{itemize}[leftmargin=1.6em,itemsep=0.15ex]
  \item (b)$\Rightarrow$(b')$\Rightarrow$(a)$\Rightarrow$(a')$\Rightarrow$(b)；
  \item (a'')$\Rightarrow$(a) 与 (b)$\Rightarrow$(a'')；
  \item (a')$\Rightarrow$(c)$\Rightarrow$(c')$\Rightarrow$(a)。
\end{itemize}
其中 (a'')$\Rightarrow$(a)、(b)$\Rightarrow$(b') 与 (c)$\Rightarrow$(c') 显然；(b')$\Rightarrow$(a) 与 (b)$\Rightarrow$(a'') 由习题 \ref{exc:28.2.2} 得到。

为证 (a)$\Rightarrow$(a')，设 $x\in X$，$U\subseteq X$ 是 $x$ 的仿射邻域且 $\cL|_U\simeq\cO_U$。由 (a)，存在齐次元素
\[
f\in S_n=\Gamma(X,\cO(n)),\quad n>0,
\]
使 $X_f\subseteq U$。于是 $X_f=U_{f|_U}$ 是截面
\[
f|_U\in\Gamma(U,\cL|_U^{\otimes n})\simeq\Gamma(U,\cO_U)
\]
的主开子集。$U$ 仿射，故 $U_{f|_U}$ 也仿射。因此 $X_f$ 是含 $x$ 的仿射主开子集，即 (a') 成立。

为证 (a')$\Rightarrow$(b)，设 $f\in S_+$ 是齐次元素且 $X_f$ 仿射。由习题 \ref{exc:28.2.2} 有
\[
\mathfrak r_{\cL}^{-1}\bigl(U_+(f)\bigr)=X_f.
\]
断言所得态射 $X_f\to U_+(f)$ 是同构。注意这一断言确实蕴含 $\mathfrak r_{\cL}$ 是像稠密的开浸入\footnote{稠密性来自「空概形是仿射的」这一观察。}。

为证此断言，先证 $X$ 拟分离。设 $X=\bigcup_{i\in I}X_{f_i}$ 是 $X$ 的有限覆盖，其中每个 $X_{f_i}$ 仿射。利用 $X$ 拟紧的假设，容易看出任意主开子集 $X_g$ 都拟紧。特别地，对每对 $(i,j)\in I^2$，交
\[
X_{f_i}\cap X_{f_j}=X_{f_if_j}
\]
拟紧。这说明 $X$ 拟分离。

由于 $X$ 拟紧且拟分离，由定理 \ref{thm:28.1.5}，典范同构
\[
S_{(f_i)}\xrightarrow{\simeq}\Gamma(X_{f_i},\cO_X)
\]
是同构。由构造可见，这给出所需的同构 $X_f\to U_+(f)$。

为证 (a')$\Rightarrow$(c)，设
\[
X=\bigcup_{i\in I}X_{f_i}
\]
是开覆盖，其中 $f_i\in S_{d_i}$ 且 $X_{f_i}$ 仿射。只需证对每个 $i\in I$，
\[
\bigoplus_{n\ge N}\cO(-n)|_{X_{f_i}}\otimes_{\Z}\Gamma(X,\cF(n)) \to \cF|_{X_{f_i}}
\]
是满射。$X_{f_i}$ 仿射，故只需证任意截面 $u\in\cF(X_{f_i})$ 落在像中。而这由注记 \ref{rem:28.1.6} 的断言 (2) 得到。

为证 (c')$\Rightarrow$(a)，设 $x\in X$，$U\subseteq X$ 是 $x$ 的仿射邻域。如下构造含 $x$ 的主开子集 $X_f\subseteq U$。设 $\cI\subseteq\cO_X$ 是对应于闭子概形 $(X\setminus U)_{\mathrm{red}}$ 的拟凝聚理想。注意 $\cI|_U\simeq\cO_U$。由 (c')，存在 $n\in\Z$ 与
\[
f\in\Gamma(X,\cI(n))\subseteq\Gamma(X,\cO(n))=S_n,
\]
使态射
\[
f\colon\cO\to\cI(n)
\]
在 $x$ 处非零。把 $f$ 看作线丛 $\cO(n)$ 的截面，并考虑相应的主开子集 $X_f$。上述非零性蕴含
\[
x\in U_{f|_U}\subseteq X_f.
\]
另一方面，对任意点 $y\in X\setminus U$，有
\[
f_y\in\cI(n)_y\simeq\cI_y\otimes_{\cO_{X,y}}\cL_y^{\otimes n}\subseteq\mathfrak m_y\otimes_{\cO_{X,y}}\cL_y^{\otimes n}.
\]
故 $y\notin X_f$，从而 $X_f\subseteq U$，如所欲证。
\end{proof}

\begin{lemma}\label{lem:28.3.2}
设 $X$ 是拟紧概形且存在丰沛线丛，则 $X$ 分离。
\end{lemma}

\begin{proof}
由 $X$ 是 $\Proj(S)$ 的开子概形，而后者分离。
\end{proof}

\begin{proposition}\label{prop:28.3.3}
设 $X$ 是拟紧且拟分离的概形，$\cL$ 是 $X$ 上的线丛。以下条件等价：
\begin{enumerate}[label=(\alph*),leftmargin=2.6em,itemsep=0.3ex]
  \item 线丛 $\cL$ 丰沛。
  \item[\textup{(d)}] 对任意有限型的拟凝聚 $\cO_X$-模 $\cF$，当 $n\gg 0$ 时 $\cF\otimes_{\cO_X}\cL^{\otimes n}$ 由有限多个整体截面生成。
  \item[\textup{(d')}] 对任意有限型的拟凝聚理想 $\cI\subseteq\cO_X$，存在 $n>0$ 使 $\cI\otimes_{\cO_X}\cL^{\otimes n}$ 由整体截面生成。
\end{enumerate}
\end{proposition}

\begin{yizhu}
(d) 是实际验证丰沛性最常用的形式：它只涉及有限型的 $\cF$ 与「足够大的扭转由有限多个截面生成」，不必逐点找仿射主开集。Stacks 把这条列为等价条件 (7)，但先要求 $X$ 拟分离，因为要用「拟凝聚模是有限型子模的滤过余极限」把一般 $\cF$ 化到有限型。
[来源：Stacks Project, Tag 01Q3（Proposition 28.27.13，条件 (7)：$X$ 拟分离且对有限型 $\cF$，$\cF\otimes\cL^{\otimes n}$ 对 $n\gg0$ 整体生成），\url{https://stacks.math.columbia.edu/tag/01Q3}（2026-11-17 访问）。]
\end{yizhu}

\begin{proof}
沿用命题—定义 \ref{prop:28.3.1} 证明中的记号。(d)$\Rightarrow$(d') 显然。

为证 (d') 蕴含 $\cL$ 丰沛，验证命题—定义 \ref{prop:28.3.1} 的条件 (c')。设 $\cI\subseteq\cO$ 是拟凝聚理想。由定理 \ref{thm:18.6.6}，$\cI$ 可写成滤过余极限 $\cI\simeq\colim\cI_\alpha$，其中每个 $\cI_\alpha$ 是有限型拟凝聚理想。由 (d')，
\[
\bigoplus_{n>0}\cO(-n)\otimes_{\Z}\Gamma(X,\cI(n)) \to \cI
\]
的像包含每个 $\cI_\alpha$。故它是满射，如所欲证。

剩下要证：若 $\cL$ 丰沛，则 (d) 成立。因 $\cL$ 丰沛，存在有限覆盖
\[
X=\bigcup_{i\in I}X_{f_i},
\]
其中 $f_i\in S_+$ 是齐次元素且 $X_{f_i}$ 仿射。注意用 $f_i$ 的正幂替换 $f_i$ 不改变 $X_{f_i}$。故可设 $f_i$（$i\in I$）次数相同，等于 $d$，其中正整数 $d$ 只依赖 $X$ 与 $\cL$。

设 $\cF$ 是有限型拟凝聚 $\cO$-模。对每个 $i\in I$，因 $X_{f_i}$ 仿射，限制 $\cF|_{X_{f_i}}$ 由有限多个截面
\[
u_{ij}\in\Gamma(X_{f_i},\cF),\quad j\in J_i
\]
生成。由注记 \ref{rem:28.1.6}(2)，当 $k\gg 0$ 时每个截面
\[
u_{ij}\otimes (f_i|_{X_{f_i}})^{\otimes k}\in\Gamma(X_{f_i},\cF(kd))
\]
可延拓为截面
\[
s_{ij}\in\Gamma(X,\cF(kd)).
\]
由构造，$\cF(kd)$ 由这些截面生成。

于是，对固定的 $\cF$，当 $k\gg 0$ 时 $\cF(kd)$ 由有限多个整体截面生成。

把此结论依次施于 $\cF(0),\cF(1),\ldots,\cF(d-1)$，即知：当 $n\gg 0$ 时 $\cF(n)$ 由有限多个整体截面生成，如所欲证。
\end{proof}

\begin{example}\label{ex:28.3.4}
拟仿射概形上的任意线丛都丰沛，因为它满足命题—定义 \ref{prop:28.3.1} 的条件 (b')。

反之，若 $X$ 是拟紧概形且 $\cO_X$ 丰沛，则 $X$ 拟仿射（命题 \ref{prop:20.6.2}）。
\end{example}

\begin{example}\label{ex:28.3.5}
设 $S$ 是 $\Z_{\ge 0}$-分次交换环，且 $S_+$ 由 $S_1$ 中有限多个元素生成。则重言线丛 $\cO_{\Proj(S)}(1)$ 丰沛。见命题 \ref{prop:26.4.8}。
\end{example}

\begin{exercise}\label{exc:28.3.6}
设 $X$ 是拟紧概形，$\cL$ 是 $X$ 上的线丛，$d$ 是正整数。证明：$\cL$ 丰沛当且仅当 $\cL^{\otimes d}$ 丰沛。
\end{exercise}

\begin{exercise}\label{exc:28.3.7}
设 $X$ 是拟紧概形，$\cL_1,\cL_2$ 是 $X$ 上的丰沛线丛。证明 $\cL_1\otimes\cL_2$ 丰沛。
\end{exercise}

\begin{exercise}\label{exc:28.3.8}
设 $n>0$，$d$ 是整数。证明：$\cO_{\PP^n_{\Z}}(d)$ 丰沛当且仅当 $d>0$。
\end{exercise}

\begin{warning}\label{warn:28.3.9}
设 $X$ 是拟紧概形，$\cL$ 是线丛。分次截面环 $S:=\bigoplus_{n\ge 0}\Gamma(X,\cL^{\otimes n})$ 一般是病态的。
\begin{itemize}[leftmargin=1.6em,itemsep=0.2ex]
  \item $S_1$ 可以是零\footnote{设 $X$ 是 $\C$ 上的完全超椭圆曲线，$x_1,x_2$ 是 $X$ 的 $\C$-点且 $\sigma(x_1)\neq x_2$，其中 $\sigma$ 是超椭圆对合。则对任意 $\C$-点 $y$，线丛 $\cO_X(x_1+x_2-y)$ 丰沛，但没有非零整体截面。}。
  \item $\cO_{\Proj(S)}(1)$ 可以不是可逆层\footnote{设 $\overline X:=\Spec(\C[x,y,z]/(xy-z^2))$，$X\subseteq\overline X$ 是去掉极大理想 $(x,y,z)$ 所给闭点后的补开子集。设 $\cI\subset\cO_{\overline X}$ 是对应于 $(x,z)$ 的理想层。则 $\cI|_X$ 是 $X$ 上的丰沛线丛。有 $\Proj(S)\simeq\overline X$。经由该同构，$\cO_{\Proj(S)}(1)$ 对应于 $\cO_X$-模 $\cI$，而 $\cI$ 不可逆。}。
\end{itemize}
\end{warning}

\begin{yizhu}
脚注中的第一个例子说明 $S=\bigoplus_n\Gamma(X,\cL^{\otimes n})$ 可以完全没有一次项，故不能指望它由 $S_1$ 生成；第二个例子说明 $\cO_{\Proj(S)}(1)$ 甚至可以不可逆。Stacks 仍记 $S=\Gamma_*(X,\cL)$，但把判据放在「$X_s$ 构成基且 $X\to\Proj(S)$ 是开浸入」这一侧，从而绕开 $S$ 本身的病态。
[来源：Stacks Project, Tag 01Q3（Proposition 28.27.13，条件 (1)–(3)：$S=\Gamma_*(X,\cL)$ 时 $X_s$ 成基与 $X\to\Proj(S)$ 开浸入的等价），\url{https://stacks.math.columbia.edu/tag/01Q3}（2026-11-17 访问）。]
\end{yizhu}

\subsection{相对版本}\label{sec:28.4}

本小节研究构造 \ref{con:28.2.3} 中的态射 $\mathfrak s_{\cL/Y}$。

\begin{definition}\label{def:28.4.1}
设 $f\colon X\to Y$ 是拟紧态射，$\cL$ 是 $X$ 上的线丛。若对任意仿射开子概形 $U\subseteq Y$，限制 $\cL|_{X_U}$ 都是 $X_U:=X\times_Y U$ 上的丰沛线丛，则称 $\cL$ 是 \emph{$f$-丰沛的}\footnote{其它术语：对 $f$ 相对丰沛（relatively ample）；$X$ 在 $Y$ 上丰沛（ample on $X$ over $Y$）。}。
\end{definition}

\begin{yizhu}
定义要求 $f$ 拟紧，这一条不能省：Stacks 的 $f$-ample 定义同样内置「$f$ 拟紧」，并注明这与 EGA II 4.6.1 一致。少了它，$f_*(\cL^{\otimes n})$ 未必拟凝聚，28.2.1 的 $\mathfrak r_{\cL/Y}$ 就无从谈起。紧随其后的推论 \ref{cor:28.4.2} 说明相对丰沛性还迫使 $f$ 分离，这一点在 §29 判定射影性时反复用到。
[来源：Stacks Project, Tag 01VH（Definition 29.38.1，relatively ample：$f$ 拟紧且每个仿射开 $V\subseteq S$ 上 $\cL|_{f^{-1}(V)}$ ample；EGA II 4.6.1），\url{https://stacks.math.columbia.edu/tag/01VH}（2026-11-17 访问）。]
\end{yizhu}

\begin{corollary}\label{cor:28.4.2}
设 $f\colon X\to Y$ 是拟紧态射且存在 $f$-丰沛线丛，则 $f$ 分离。
\end{corollary}

\begin{proof}
由引理 \ref{lem:28.3.2} 得到。
\end{proof}

\begin{corollary}\label{cor:28.4.3}
设 $f\colon X\to Y$ 是拟紧态射，$\cL$ 是 $X$ 上的线丛。以下条件等价：
\begin{enumerate}[label=(\alph*),leftmargin=2.6em,itemsep=0.3ex]
  \item 线丛 $\cL$ 是 $f$-丰沛的。
  \item 态射 $\mathfrak r_{\cL/Y}\colon X\to\Proj_Y(\mathcal S)$ 有定义\footnote{换言之，$f_*(\cL^{\otimes n})$ 对 $n\ge 0$ 拟凝聚，且 $U_{\cL/Y}=X$。}，且是像稠密的开浸入。
  \item[\textup{(b')}] 态射 $\mathfrak r_{\cL/Y}\colon X\to\Proj_Y(\mathcal S)$ 有定义，且诱导到其像上的同胚。
\end{enumerate}
若 $X$ 拟分离，则上述条件还等价于：
\begin{enumerate}[label=(\alph*),leftmargin=2.6em,itemsep=0.3ex,start=3]
  \item 对任意拟凝聚 $\cO_X$-模 $\cF$，典范态射
  \[
  \bigoplus_{n\ge N}\Bigl(f_*f^*\bigl(\cF\otimes_{\cO_X}\cL^{\otimes n}\bigr)\Bigr)\otimes_{\cO_X}\cL^{\otimes -n} \to \cF
  \]
  对任意整数 $N$ 都是满射。
  \item[\textup{(c')}] 对任意拟凝聚理想 $\cI\subseteq\cO_X$，典范态射
  \[
  \bigoplus_{n\ge N}\Bigl(f_*f^*\bigl(\cI\otimes_{\cO_X}\cL^{\otimes n}\bigr)\Bigr)\otimes_{\cO_X}\cL^{\otimes -n} \to \cI
  \]
  都是满射。
\end{enumerate}
若 $X$ 拟紧且拟分离，则上述条件还等价于：
\begin{enumerate}[label=(\alph*),leftmargin=2.6em,itemsep=0.3ex,start=5]
  \item[\textup{(d)}] 对任意有限型拟凝聚 $\cO_X$-模 $\cF$，存在 $N$，使得对任意 $n\ge N$ 都存在有限型子模
  \[
  \cM\subseteq f_*\bigl(\cF\otimes_{\cO_X}\cL^{\otimes n}\bigr),
  \]
  使
  \[
  f^*\cM \to f^*f_*\bigl(\cF\otimes_{\cO_X}\cL^{\otimes n}\bigr) \to \cF\otimes_{\cO_X}\cL^{\otimes n}
  \]
  是满射。
  \item[\textup{(d')}] 对任意有限型拟凝聚理想 $\cI\subseteq\cO_X$，存在 $n>0$ 使
  \[
  f^*f_*\bigl(\cI\otimes_{\cO_X}\cL^{\otimes n}\bigr) \to \cI\otimes_{\cO_X}\cL^{\otimes n}
  \]
  是满射。
\end{enumerate}
\end{corollary}

\begin{proofsketch}
(a)$\Leftrightarrow$(b)$\Leftrightarrow$(b') 由命题—定义 \ref{prop:28.3.1} 得到。

当 $X$ 拟分离时，注意对任意仿射开子概形 $U\subseteq Y$，开浸入 $j\colon X_U\to X$ 拟紧且拟分离，故 $j_*$ 保持拟凝聚模。于是可以应用命题—定义 \ref{prop:28.3.1} 得到 (a)$\Leftrightarrow$(c)$\Leftrightarrow$(c')。

当 $X$ 拟紧且拟分离时，可应用命题 \ref{prop:28.3.3}、定理 \ref{thm:18.6.1} 与定理 \ref{thm:18.6.6} 得到 (a)$\Leftrightarrow$(d)$\Leftrightarrow$(d')。
\end{proofsketch}

\begin{lemma}\label{lem:28.4.4}
相对丰沛线丛有以下标准性质。
\begin{enumerate}[label=(\arabic*),leftmargin=2.6em,itemsep=0.35ex]
  \item 设 $f\colon X\to Y$ 是拟仿射态射，则 $X$ 上的任意线丛 $\cL$ 都是 $f$-丰沛的。
  \item 设 $f\colon X\to Y$ 拟紧，$d>0$ 是正整数。则 $\cL$ 是 $f$-丰沛的当且仅当 $\cL^{\otimes d}$ 是。
  \item 设 $f\colon X\to Y$ 拟紧，$\cL$（相应地 $\cK$）是 $X$（相应地 $Y$）上的线丛。则 $\cL$ 是 $f$-丰沛的当且仅当 $\cL\otimes_{\cO_X}f^*\cK$ 是。
  \item 设 $X\xrightarrow{f}Y\xrightarrow{g}Z$ 是拟紧态射，且 $f$ 诱导到其像上的同胚。设 $\cK$ 是 $g$-丰沛线丛，则 $f^*\cK$ 是 $(g\circ f)$-丰沛的。
  \item[\textup{(4')}] 设 $X\xrightarrow{f}Y\xrightarrow{g}Z$ 是拟紧态射，且 $Z$ 拟紧。设 $\cL$ 是 $f$-丰沛线丛，$\cK$ 是 $g$-丰沛线丛。则当 $l\gg 0$ 时 $\cL\otimes_{\cO_X}f^*(\cK^{\otimes l})$ 是 $(g\circ f)$-丰沛的。
  \item 设 $X\xrightarrow{f}Y\xrightarrow{g}Z$ 是概形态射，且 $g\circ f$ 拟紧、$g$ 拟分离。设 $\cL$ 是 $(g\circ f)$-丰沛线丛，则 $\cL$ 是 $f$-丰沛线丛。
  \item 相对丰沛线丛对张量积稳定。
  \item 相对丰沛线丛经基变换保持。
  \item 相对丰沛性对目标概形是局部的。
\end{enumerate}
\end{lemma}

\begin{proof}
(1) 由例 \ref{ex:28.3.4} 得到。(2) 由习题 \ref{exc:28.3.6} 得到。(5) 通过验证推论 \ref{cor:28.4.3} 的条件 (d') 得到。(6) 由习题 \ref{exc:28.3.7} 得到。(7) 通过验证推论 \ref{cor:28.4.3} 的条件 (b') 得到。(8) 由定义显然。(3) 由 (8) 得到，因为可以在 $Y$ 上局部平凡化 $\cK$。

为证 (4)，可设 $Z$ 仿射且 $\cK$ 丰沛。此时验证命题—定义 \ref{prop:28.3.1} 的条件 (a) 即知 $f^*\cK$ 丰沛。

剩下证 (4')。由 (8)，可设 $Z$ 仿射且 $\cK$ 丰沛。与命题 \ref{prop:28.3.3} 的证明一样，可以取 $d>0$ 与有限多个截面
\[
v_i\in\Gamma(Y,\cK^{\otimes d}),\quad i\in I,
\]
使
\[
Y=\bigcup_{i\in I}Y_{v_i}
\]
且每个 $Y_{v_i}$ 仿射。注意 $d$ 只依赖 $Y$ 与 $\cK$。

令
\[
u_i\in\Gamma(X,f^*(\cK)^{\otimes d})
\]
是截面 $v_i$ 的拉回。则
\[
X_{u_i}\simeq X\times_Y Y_{v_i},
\]
从而
\[
X=\bigcup_{i\in I}X_{u_i}.
\]
又按定义，$\cL|_{X_{u_i}}$ 是 $X_{u_i}$ 上的丰沛线丛。故如前可取 $n>0$ 与有限多个截面
\[
w_{ij}\in\Gamma(X_{u_i},\cL^{\otimes n}),\quad j\in J_i,
\]
使
\[
X_{u_i}=\bigcup_{j\in J_i}(X_{u_i})_{w_{ij}}
\]
且每个 $(X_{u_i})_{w_{ij}}$ 仿射。注意 $X$ 拟紧且分离，因为 $f$ 与 $g$ 都如此。故由定理 \ref{thm:28.1.5}，当 $m\gg 0$ 时截面
\[
w_{ij}\otimes (u_i|_{X_{u_i}})^{\otimes m}\in\Gamma\bigl(X_{u_i},\cL^{\otimes n}\otimes_{\cO_X}f^*(\cK)^{\otimes md}\bigr)
\]
可延拓为截面
\[
s_{ij}\in\Gamma\bigl(X,\cL^{\otimes n}\otimes_{\cO_X}f^*(\cK)^{\otimes md}\bigr).
\]
由构造有
\[
X_{s_{ij}}=(X_{u_i})_{w_{ij}}.
\]
于是
\[
X=\bigcup_{i\in I,\,j\in J_i}X_{s_{ij}}
\]
是 $X$ 的覆盖，且每个 $X_{s_{ij}}$ 仿射。取 $m=nr$，即知
\[
\cL\otimes_{\cO_X}f^*(\cK)^{\otimes rd}
\]
当 $r\gg 0$ 时丰沛，因为它满足命题—定义 \ref{prop:28.3.1} 的条件 (a')。

至此证明了：
\begin{itemize}[leftmargin=1.6em,itemsep=0.2ex]
  \item[(*)] 对任意 $f$-丰沛线丛 $\cL$，线丛 $\cL\otimes_{\cO_X}f^*(\cK)^{\otimes rd}$ 当 $r\gg 0$ 时丰沛。
\end{itemize}
注意由 (3)，对任意整数 $e$，$\cL\otimes_{\cO_X}f^*(\cK)^{\otimes e}$ 都是 $f$-丰沛的。故可把 (*) 施于 $\cL\otimes_{\cO_X}f^*(\cK)^{\otimes e}$（$e\in\{0,1,\cdots,d-1\}$），得知 $\cL\otimes_{\cO_X}f^*(\cK)^{\otimes l}$ 当 $l\gg 0$ 时丰沛。
\end{proof}

\begin{warning}\label{warn:28.4.5}
在引理 \ref{lem:28.4.4}(4') 中，$\cL\otimes_{\cO_X}f^*\cK$ 一般不是 $(g\circ f)$-丰沛的。事实上，设这总成立，则
\[
\cL\simeq\bigl(\cL\otimes_{\cO_X}f^*(\cK^{-1})\bigr)\otimes_{\cO_X}f^*(\cK)
\]
会 $(g\circ f)$-丰沛，因为由引理 \ref{lem:28.4.4}(3)，$\cL\otimes_{\cO_X}f^*(\cK^{-1})$ 也是 $f$-丰沛的。换言之，我们将得到「任意 $f$-丰沛线丛都是 $(g\circ f)$-丰沛的」，而取 $f=\id$ 时这显然荒谬。
\end{warning}

\begin{corollary}\label{cor:28.4.6}
设 $Y$ 是概形，$\mathcal S$ 是 $\Z_{\ge 0}$-分次拟凝聚 $\cO_Y$-代数，且
\begin{itemize}[leftmargin=1.6em,itemsep=0.2ex]
  \item $\mathcal S_+$ 由 $\mathcal S_1$ 生成；
  \item $\mathcal S_1$ 有限型。
\end{itemize}
设 $i\colon X\to\Proj_Y(\mathcal S)$ 是拟紧局部闭浸入。则 $i^*(\cO(1))$ 对态射 $X\to Y$ 相对丰沛。
\end{corollary}

\begin{proof}
把引理 \ref{lem:28.4.4}(4) 施于 $X\to\Proj_Y(\mathcal S)\to Y$ 即得。
\end{proof}

\subsection{极丰沛线丛}\label{sec:28.5}

本小节研究构造 \ref{con:28.2.1} 中的态射 $\mathfrak r_{\cL}$。

\begin{definition}\label{def:28.5.1}
设 $f\colon X\to Y$ 是概形态射，$\cL$ 是 $X$ 上的线丛。若存在拟凝聚 $\cO_Y$-模 $\cE$ 与在 $Y$ 上定义的局部闭浸入
\[
i\colon X\to\PP_Y(\cE),
\]
使 $\cL\simeq i^*(\cO_{\PP_Y(\cE)}(1))$，则称 $\cL$ 是 \emph{$f$-极丰沛的}。

若 $X$ 上的线丛 $\cL$ 对态射 $X\to\Spec(\Z)$ 相对极丰沛，则称 $\cL$ 是\emph{极丰沛的}。
\end{definition}

\begin{yizhu}
Stacks 的定义与此逐字相同：存在拟凝聚 $\cO_S$-模 $\cE$ 与 $S$-浸入 $i\colon X\to\mathbf P(\cE)$，使 $\cL\cong i^*\cO(1)$；其中 immersion 指局部闭浸入，且同样不要求 $f$ 拟紧。这与下面的警告 \ref{warn:28.5.2} 一致：拟紧性来自历史惯例，而非定义本身。定义并不要求 $\cE$ 有限型；当 $Y$ 拟紧且 $f$ 有限型时，命题 \ref{prop:28.5.6} 说明可取 $\cE\subseteq f_*\cL$ 为有限型。
[来源：Stacks Project, Tag 01VM（Definition 29.39.1，very ample invertible sheaf），\url{https://stacks.math.columbia.edu/tag/01VM}（2026-11-17 访问）。]
\end{yizhu}

\begin{warning}\label{warn:28.5.2}
出于历史原因，在上述定义中人们不假设 $f$ 拟紧\footnote{不过，笔者从未见过这一概念在非拟紧情形下的任何应用。}。
\end{warning}

与前面一样，有：

\begin{lemma}\label{lem:28.5.3}
设 $f\colon X\to Y$ 是存在 $f$-极丰沛线丛的态射，则 $f$ 分离。
\end{lemma}

下述结果表明，在拟紧情形下，定义 \ref{def:28.5.1} 中的 $\cO_Y$-模 $\cE$ 有典范选法。

\begin{lemma}\label{lem:28.5.4}
设 $f\colon X\to Y$ 是概形间的拟紧态射。则 $X$ 上的线丛 $\cL$ 是 $f$-极丰沛的，当且仅当它满足以下条件：
\begin{enumerate}[label=(\roman*),leftmargin=2.6em,itemsep=0.3ex]
  \item $\cO_Y$-模 $f_*(\cL)$ 拟凝聚，且典范态射
  \[
  v\colon f^*f_*(\cL)\to\cL
  \]
  是满射；
  \item 典范态射（见习题 \ref{exc:28.2.5}）
  \[
  \mathfrak s_{\cL/Y}\colon X\to\PP_Y(f_*\cL)
  \]
  是局部闭浸入。
\end{enumerate}
\end{lemma}

\begin{proof}
「若」由同构
\[
\mathfrak s_{\cL/Y}^*\bigl(\cO_{\PP_Y(f_*\cL)}(1)\bigr)\simeq\cO_{\PP_X(\cL)}(1)\simeq\cL
\]
得到。

为证「仅当」，设 $i\colon X\to\PP_Y(\cE)$ 是定义 \ref{def:28.5.1} 中的局部闭浸入。注意 $f_*(\cL)$ 拟凝聚，因为 $f$ 拟紧且分离。由推论 \ref{cor:27.5.6}，态射 $i$ 由满射 $\psi\colon f^*\cE\to\cL$ 分类。该态射分解为
\[
f^*\cE\xrightarrow{f^*(\varphi_1)}f^*f_*\cL\xrightarrow{v}\cL,
\]
其中 $\varphi_1\colon\cE\to f_*\cL$ 是对应于 $\psi$ 的态射。这说明 $v\colon f^*f_*\cL\to\cL$ 也是满射。此即 (i)。

为证 (ii)，考虑由 $\varphi_1$ 诱导的分次同态
\[
\varphi\colon\Sym_{\cO_Y}(\cE)\to\Sym_{\cO_Y}(f_*\cL).
\]
它诱导 $Y$-态射
\[
g\colon U_+(\varphi)\to\PP_Y(\cE),
\]
其中 $U_+(\varphi)$ 是 $\PP_Y(f_*\cL)$ 的开子集。展开定义即知 $\mathfrak s_{\cL/Y}$ 的像含于 $U_+(\varphi)$，且
\[
g\circ\mathfrak s_{\cL/Y}=i.
\]
因 $i$ 是局部闭浸入，$\mathfrak s_{\cL/Y}$ 也是。此即 (ii)。
\end{proof}

\begin{corollary}\label{cor:28.5.5}
设 $f\colon X\to Y$ 是拟紧态射。对 $X$ 上的线丛而言，$f$-极丰沛性对 $Y$ 是局部的。
\end{corollary}

下述结果表明，在更强的有限性假设下，可以把定义 \ref{def:28.5.1} 中的 $\cO_Y$-模 $\cE$ 取成有限型。

\begin{proposition}\label{prop:28.5.6}
设 $Y$ 是拟紧且拟分离的概形，$f\colon X\to Y$ 是有限型态射。则 $X$ 上的线丛 $\cL$ 是 $f$-极丰沛的，当且仅当存在拟凝聚子模 $\cE\subseteq f_*\cL$，使
\begin{itemize}[leftmargin=1.6em,itemsep=0.25ex]
  \item $\cO_Y$-模 $\cE$ 有限型，且典范态射 $f^*\cE\to\cL$ 是满射；
  \item 典范态射 $X\to\PP_Y(\cE)$ 是局部闭浸入。
\end{itemize}
\end{proposition}

\begin{proof}
「若」显然。

为证「仅当」，设 $\cL$ 是 $X$ 上的 $f$-极丰沛线丛。由引理 \ref{lem:28.5.4}，典范态射 $v\colon f^*f_*(\cL)\to\cL$ 是满射，且诱导局部闭浸入
\[
r\colon X\to\PP_Y(f_*\cL).
\]
对任意 $\cE\in \Sub(f_*\cL)_{\mathrm{ft}}$，令
\[
r_{\cE}\colon X\dashrightarrow \PP_Y(\cE)
\]
为对应于复合 $f^*\cE\to f^*f_*(\cL)\to\cL$ 的部分定义态射。由定理 \ref{thm:18.6.6}，有
\[
\colim_{\cE\in \Sub(f_*\cL)_{\mathrm{ft}}}\cE\xrightarrow{\simeq}f_*\cL.
\]
因此对足够大的对象 $\cE\in \Sub(f_*\cL)_{\mathrm{ft}}$，典范态射
\[
v_{\cE}\colon f^*\cE\to\cL
\]
是满射，从而 $r_{\cE}$ 在 $X$ 上有定义。

注意对 $\Sub(f_*\cL)_{\mathrm{ft}}$ 中足够大的对象 $\cE\subseteq\cE'$，有部分定义态射 $r_{\cE,\cE'}\colon\PP_Y(\cE')\dashrightarrow\PP_Y(\cE)$，使 $r_{\cE'}$ 的像含于 $\PP_Y(\cE')$ 中 $r_{\cE,\cE'}$ 有定义的开子集。此外 $r_{\cE}=r_{\cE,\cE'}\circ r_{\cE'}$。故若 $r_{\cE}$ 是局部闭浸入，则 $r_{\cE'}$ 也是。

设 $Y=\bigcup_{i\in I}Y_i$ 是 $Y$ 由仿射开子概形组成的有限覆盖。要证 $r_{\cE}$ 是局部闭浸入，只需证它到每个 $Y_i$ 的基变换是。故由上一段，只需证对每个 $i\in I$，$r_{\cE}$ 到 $Y_i$ 的基变换对足够大的 $\cE$ 是局部闭浸入。这使我们可以化到 $Y=\Spec(R)$ 仿射的情形。

设
\[
s\in\Gamma\bigl(Y,\Sym^d_{\cO_Y}(f_*\cL)\bigr),\quad d>0,
\]
是整体截面，考虑相应的标准开子集 $U_+(s)\subseteq\PP_Y(f_*\cL)$。注意对任意正整数 $n$ 有 $U_+(s)=U_+(s^n)$。由构造有
\[
r^{-1}\bigl(U_+(s)\bigr)=X_{v(s)},
\]
其中 $v(s)$ 是 $s$ 在复合
\[
\Gamma\bigl(Y,\Sym^d_{\cO_Y}(f_*\cL)\bigr)\to\Gamma\bigl(X,f^*\Sym^d_{\cO_Y}(f_*\cL)\bigr)\xrightarrow{v^{\otimes d}}\Gamma(X,\cL^{\otimes d})
\]
下的像。因 $r$ 是局部闭浸入且 $X$ 拟紧，存在 $d>0$ 与有限多个截面
\[
s_i\in\Gamma\bigl(Y,\Sym^d_{\cO_Y}(f_*\cL)\bigr),\quad i\in I,
\]
使 $X=\bigcup_{i\in I}X_{v(s_i)}$，且每个态射
\[
X_{v(s_i)}\to U_+(s_i)
\]
都是闭浸入。特别地，每个 $X_{v(s_i)}$ 仿射，且有满射
\[
\Gamma\bigl(Y,\Sym_{\cO_Y}(f_*\cL)\bigr)_{(s_i)}\to\Gamma(X_{v(s_i)},\cO_X).
\]

注意每个 $X_{v(s_i)}$ 在 $Y$ 上有限型，因为 $X$ 如此。故 $\Gamma(X_{v(s_i)},\cO_X)$ 是有限生成 $A$-代数。于是可取足够大的对象 $\cE\in\Sub(f_*\cL)_{\mathrm{ft}}$，使：
\begin{itemize}[leftmargin=1.6em,itemsep=0.25ex]
  \item 每个 $s_i$ 含于
  \[
  \Gamma\bigl(Y,\Sym^d_{\cO_Y}(\cE)\bigr)\to\Gamma\bigl(Y,\Sym^d_{\cO_Y}(f_*\cL)\bigr)
  \]
  的像中；
  \item 设 $s'_i\in\Gamma(Y,\Sym^d_{\cO_Y}(\cE))$ 是像为 $s_i$ 的截面，则复合
  \[
  \Gamma\bigl(Y,\Sym_{\cO_Y}(\cE)\bigr)_{(s'_i)}\to\Gamma\bigl(Y,\Sym_{\cO_Y}(f_*\cL)\bigr)_{(s_i)}\to\Gamma(X_{v(s_i)},\cO_X)
  \]
  是满射。
\end{itemize}
这说明态射 $r_{\cE}\colon X\to\PP_Y(\cE)$ 限制为闭浸入
\[
X_{v(s_i)}=X_{v_{\cE}(s'_i)}\to U_+(s'_i).
\]
因 $X$ 被这些 $X_{v(s_i)}$ 覆盖，故 $r_{\cE}$ 是局部闭浸入，如所欲证。
\end{proof}

\begin{corollary}\label{cor:28.5.7}
设 $Y$ 是拟紧概形且存在丰沛线丛 $\cK$。设 $f\colon X\to Y$ 是有限型态射且存在 $f$-极丰沛线丛。则存在 $n\ge 0$ 与局部闭浸入 $i\colon X\to\PP^n_Y$。
\end{corollary}

\begin{proof}
由命题 \ref{prop:28.5.6}，有在 $Y$ 上定义的局部闭浸入 $X\to\PP_Y(\cE)$，其中 $\cE$ 是有限型拟凝聚 $\cO_Y$-模。容易看出，对 $Y$ 上任意线丛 $\cL'$，
\[
\PP_Y(\cE\otimes_{\cO_Y}\cL')\simeq\PP_Y(\cE).
\]
因 $\cK$ 丰沛，可用 $\cE\otimes_{\cO_Y}\cK^{\otimes d}$ 替换 $\cE$（$d$ 足够大），从而设 $\cE$ 由整体截面生成。取满射 $\cO_Y^{\oplus n+1}\to\cE$，它诱导闭浸入
\[
\PP_Y(\cE)\to\PP_Y(\cO_Y^{\oplus n+1})\simeq\PP^n_Y.
\]
于是复合 $X\to\PP_Y(\cE)\to\PP^n_Y$ 给出所需的局部闭浸入。
\end{proof}

\begin{lemma}\label{lem:28.5.8}
设 $f\colon X\to Y$ 是概形间的态射，$\cL$ 是 $X$ 上的 $f$-极丰沛线丛。设 $\cL'$ 是 $X$ 上的线丛，且存在 $\cE'\in\QCoh(\cO_Y)$ 与满射 $f^*\cE'\to\cL'$。则 $\cL\otimes_{\cO_X}\cL'$ 是 $f$-极丰沛的。
\end{lemma}

\begin{proof}
按假设，有在 $Y$ 上定义的局部闭浸入
\[
r\colon X\to\PP_Y(\cE)=:P,
\]
使 $r^*(\cO_P(1))\simeq\cL$，其中 $\cE$ 是某个拟凝聚 $\cO_Y$-模。又有在 $Y$ 上定义的态射
\[
r'\colon X\to\PP_Y(\cE')=:P',
\]
使 $(r')^*(\cO_{P'}(1))\simeq\cL'$。考虑复合
\[
s\colon X\xrightarrow{\Delta_f}X\times_Y X\xrightarrow{r\times r'}P\times_Y P'\to\PP_Y(\cE\otimes_{\cO_Y}\cE'),
\]
其中最后一个态射是 Segre 嵌入（习题 \ref{exc:27.6.2}）。于是
\[
s^*(\cO(1))\simeq\cL\otimes_{\cO_X}\cL'.
\]
又因
\[
X\xrightarrow{\Delta_f}X\times_Y X\xrightarrow{r\times r'}P\times_Y P'\to P
\]
等于 $r$，而 $r$ 是局部闭浸入，可见
\[
X\xrightarrow{\Delta_f}X\times_Y X\xrightarrow{r\times r'}P\times_Y P'
\]
是局部闭浸入。故 $s$ 是局部闭浸入。这说明 $\cL\otimes_{\cO_X}\cL'$ 极丰沛，如所欲证。
\end{proof}

\begin{corollary}\label{cor:28.5.9}
设 $f\colon X\to Y$ 是概形间的态射。则 $f$-极丰沛线丛对张量积稳定。
\end{corollary}

\begin{corollary}\label{cor:28.5.10}
设 $f\colon X\to Y$ 是概形间的态射，$\cL$（相应地 $\cK$）是 $X$（相应地 $Y$）上的线丛。则 $\cL$ 是 $f$-极丰沛的当且仅当 $\cL\otimes_{\cO_X}f^*\cK$ 是。
\end{corollary}

下述结果的证明留给读者。

\begin{lemma}\label{lem:28.5.11}
相对极丰沛线丛有以下标准性质。
\begin{enumerate}[label=(\arabic*),leftmargin=2.6em,itemsep=0.35ex]
  \item $X$ 上的任意线丛 $\cL$ 都是 $\id_X$-极丰沛的。
  \item 设 $X\xrightarrow{f}Y\xrightarrow{g}Z$ 是概形态射，且 $f$ 是局部闭浸入。设 $\cK$ 是 $g$-极丰沛线丛，则 $f^*\cK$ 是 $(g\circ f)$-极丰沛的。
  \item[\textup{(2')}] 设 $X\xrightarrow{f}Y\xrightarrow{g}Z$ 是拟紧态射，$Z$ 拟紧，$f$ 有限型。设 $\cL$ 是 $f$-极丰沛线丛，$\cK$ 是 $g$-极丰沛线丛。则当 $l\gg 0$ 时 $\cL\otimes_{\cO_X}f^*(\cK^{\otimes l})$ 是 $(g\circ f)$-极丰沛的。
  \item 设 $X\xrightarrow{f}Y\xrightarrow{g}Z$ 是概形态射。设 $\cL$ 是 $(g\circ f)$-极丰沛线丛，则 $\cL$ 是 $f$-极丰沛线丛。
  \item 相对极丰沛线丛经基变换保持。
\end{enumerate}
\end{lemma}

\subsection{丰沛与极丰沛}\label{sec:28.6}

\begin{lemma}\label{lem:28.6.1}
设 $f\colon X\to Y$ 是拟紧态射。任意 $f$-极丰沛线丛都是 $f$-丰沛的。
\end{lemma}

\begin{proof}
由引理 \ref{lem:28.5.4}，有有定义的局部闭浸入
\[
\mathfrak s_{\cL/Y}\colon X\to\PP_Y(f_*\cL).
\]
由习题 \ref{exc:28.2.4}，有有定义的态射
\[
\mathfrak r_{\cL/Y}\colon X\to\Proj_Y\Bigl(\bigoplus_{n\ge 0}f_*(\cL^{\otimes n})\Bigr).
\]
由构造，分次同态
\[
\Sym_{\cO_Y}(f_*\cL)\to\bigoplus_{n\ge 0}f_*(\cL^{\otimes n})
\]
诱导部分定义态射
\[
g\colon\Proj_Y\Bigl(\bigoplus_{n\ge 0}f_*(\cL^{\otimes n})\Bigr)\dashrightarrow\PP_Y(f_*\cL),
\]
使 $g$ 在 $\mathfrak r_{\cL/Y}$ 的像上有定义，且
\[
\mathfrak s_{\cL/Y}=g\circ\mathfrak r_{\cL/Y}.
\]
因 $\mathfrak s_{\cL/Y}$ 是局部闭浸入，$\mathfrak r_{\cL/Y}$ 也是。故 $\cL$ 是 $f$-丰沛的，如所欲证。
\end{proof}

\begin{proposition}\label{prop:28.6.2}
设 $Y$ 是拟紧概形，$f\colon X\to Y$ 是有限型态射，$\cL$ 是 $X$ 上的线丛。以下条件等价：
\begin{enumerate}[label=(\roman*),leftmargin=2.6em,itemsep=0.3ex]
  \item 线丛 $\cL$ 是 $f$-丰沛的。
  \item 当 $n\gg 0$ 时，线丛 $\cL^{\otimes n}$ 是 $f$-极丰沛的。
  \item 存在 $n>0$ 使 $\cL^{\otimes n}$ 是 $f$-丰沛的。
\end{enumerate}
\end{proposition}

\begin{proof}
(ii)$\Rightarrow$(iii) 显然。(iii)$\Rightarrow$(i) 由引理 \ref{lem:28.6.1} 与引理 \ref{lem:28.4.4}(2) 得到。

为证 (i)$\Rightarrow$(ii)，设 $\cL$ 是 $f$-丰沛的。由推论 \ref{cor:28.5.5}，可化到 $Y=\Spec(A)$ 仿射的情形。注意此时 $\cL$ 是绝对丰沛的。于是断言由下述更强的结果得到。
\end{proof}

\begin{lemma}\label{lem:28.6.3}
设 $f\colon X\to Y$ 是局部有限型态射。假设 $X$ 拟紧且 $\cL$ 是 $X$ 上的丰沛线丛。则对足够大的 $n$，存在局部闭浸入 $\iota\colon X\to\PP^m_Y$（其中 $m$ 依赖 $n$），使 $\iota^*(\cO_{\PP^m_Y}(1))\simeq\cL^{\otimes n}$。特别地，当 $n\gg 0$ 时 $\cL^{\otimes n}$ 是 $f$-极丰沛的。
\end{lemma}

\begin{proof}
因 $\cL$ 丰沛，由命题—定义 \ref{prop:28.3.1}(a'')，可取有限多个截面
\[
s_i\in\Gamma(X,\cL^{\otimes d_i}),\quad i\in I,
\]
使
\[
X=\bigcup_{i\in I}X_{s_i},
\]
且每个 $X_{s_i}$ 仿射，并满足 $f(X_{s_i})\subseteq Y_i$，其中 $Y_i=\Spec(A_i)$ 是 $Y$ 的某个仿射开子集。

由定理 \ref{thm:28.1.5}，有典范同构
\begin{equation}
\Bigl(\bigoplus_{n\ge 0}\Gamma(X,\cL^{\otimes n})\Bigr)_{(s_i)}\xrightarrow{\simeq}\Gamma(X_{s_i},\cO_X). \label{eq:28.2}
\end{equation}
注意 $X_{s_i}$ 在 $Y_i$ 上有限型。故 $A_i$-代数 $\Gamma(X_{s_i},\cO_X)$ 可由有限多个元素
\[
u_{ij}\in\Gamma(X_{s_i},\cO_X),\quad j\in I_j
\]
生成。令
\[
s_{ij}s_i^{-n_{ij}}\in\Bigl(\bigoplus_{n\ge 0}\Gamma(X,\cL^{\otimes n})\Bigr)_{(s_i)}
\]
为相应的元素，其中
\[
s_{ij}\in\Gamma(X,\cL^{\otimes d_in_{ij}}).
\]

取 $d_in_{ij}$ 的足够大的公倍数 $N$。考虑截面
\begin{equation}
s_i^{\otimes N/d_i},\quad s_{ij}\otimes s_i^{\otimes (N/d_i-n_{ij})}\in\Gamma(X,\cL^{\otimes N}) \label{eq:28.3}
\end{equation}
与相应的态射
\begin{equation}
f^*\bigl(\cO_Y^{\oplus(m+1)}\bigr)\simeq\cO_X^{\oplus(m+1)}\to\cL^{\otimes N}, \label{eq:28.4}
\end{equation}
其中 $m+1=|I|+\sum_{i\in I}|I_j|$。为便于引用，记
\[
T_i,\ T_{ij}\in\Gamma\bigl(Y,\cO_Y^{\oplus(m+1)}\bigr)
\]
为被态射 \eqref{eq:28.4} 送到 \eqref{eq:28.3} 的截面。

断言 \eqref{eq:28.4} 是满射。为此只需证它限制到每个 $X_{s_i}$ 上都是满射。但这显然，因为截面 $s_i^{\otimes N/d_i}$ 在同构
\[
\Gamma(X_{s_i},\cL^{\otimes n})\xrightarrow{\simeq}\Gamma(X_{s_i},\cO_X)
\]
下对应于单位截面。

因 \eqref{eq:28.4} 是满射，得到态射
\begin{equation}
\iota\colon X\to\PP^m_Y\simeq\Proj_Y\bigl(\cO_Y[\{T_i\},\{T_{ij}\}]\bigr) \label{eq:28.5}
\end{equation}
使 $\iota^*(\cO_{\PP^m_Y}(1))\simeq\cL^{\otimes N}$。断言 $\iota$ 是局部闭浸入。

由构造有
\[
\iota^{-1}\bigl(U_+(T_i)\bigr)=X_{s_i^{\otimes N/d_i}}=X_{s_i}.
\]
故只需证对任意 $k\in I$，态射 $\iota$ 诱导局部闭浸入 $X_{s_k}\to U_+(T_k)$，其中 $U_+(T_k)\subseteq\PP^m_Y$ 是标准开子集。注意该态射分解为
\[
X_{s_k}\to U_+(T_k)\times_Y Y_k\to U_+(T_k),
\]
故只需证第一个态射是闭浸入。注意 $X_{s_k}$ 与 $U_+(T_k)\times_Y Y_k$ 都仿射，而上述态射由 $A_k$-同态
\begin{equation}
A_k[\{T_i\},\{T_{ij}\}]_{(T_k)}\to\Gamma(X_{s_k},\cO_X) \label{eq:28.6}
\end{equation}
给出，它把生成元 $\{T_i\},\{T_{ij}\}$ 送到 \eqref{eq:28.4} 在态射 \eqref{eq:28.2} 下的像。特别地，每个 $u_{kj}$（$j\in J_k$）都含于像中。故 \eqref{eq:28.6} 是满射，如所欲证。这说明 \eqref{eq:28.5} 是局部闭浸入。

至此证明了引理的断言对某个 $n=N$ 成立。为证它对 $n\gg 0$ 成立，借用引理 \ref{lem:28.5.8} 证明中的技巧。即：因 $\cL$ 丰沛，当 $n'\gg 0$ 时 $\cL^{\otimes n'}$ 由整体截面生成。取这样的满射
\[
f^*\bigl(\cO_Y^{m'}\bigr)\simeq\cO_X^{m'}\to\cL^{\otimes n'},
\]
并考虑相应的态射
\[
\iota'\colon X\to\PP^{m'}_Y,
\]
使 $(\iota')^*(\cO_{\PP^{m'}_Y}(1))\simeq\cL^{\otimes n'}$。与 loc.\ cit.\ 一样，复合
\[
\jmath\colon X\to\PP^m_Y\times_Y\PP^{m'}_Y\to\PP^{mm'+m+m'}_Y
\]
是局部闭浸入，且
\[
\jmath^*(\cO(1))\simeq\cL^{\otimes(N+n')}.
\]
\end{proof}
